{
    \sectionnonum{附录}

    \appendixsubsecmajornumbering

    \subsection{核心超参数与系统配置 (Hyperparameters and Configurations)}

    \begin{table}[!htbp]
    \centering
    \caption{RAG系统关键超参数与实验配置}
    \label{tab:hyperparameters}
    \small
    \begin{tabularx}{\linewidth}{|l|c|X|}
    \hline
    \textbf{参数名称 (Parameter)} & \textbf{取值 (Value)} & \textbf{说明 (Description)} \\
    \hline
    Top-$k$ (Initial Retrieval) & 5 & 初次稠密检索召回的文档数量 \\
    \hline
    PRF Threshold ($\tau$) & 0.85 & 触发伪相关反馈（二次检索）的重排序得分阈值 \\
    \hline
    RRF $k$ (Smoothing) & 60 & 互惠排序融合算法中的平滑常数 \\
    \hline
    CoVe Threshold (Standard) & 0.50 & 验证模块判断声明成立的软置信度阈值 \\
    \hline
    CoVe Threshold (Strict) & 0.90 & 严格验证配置下的置信度阈值 \\
    \hline
    LLM Context Window & 2048 & 生成回答时使用的最大上下文词元数 \\
    \hline
    \end{tabularx}
    \end{table}

    \subsection{二次检索触发策略核心伪代码 (Pseudo-code for Adaptive PRF)}

    \begin{algorithm}[!htbp]
    \renewcommand{\algorithmicrequire}{\textbf{输入:}}
    \renewcommand{\algorithmicensure}{\textbf{输出:}}
    \caption{基于置信度的二次检索触发与融合策略}
    \begin{algorithmic}[1]
    \REQUIRE 查询 $q$，语料库 $\mathcal{D}$，初始召回数 $k$，置信度阈值 $\tau$
    \ENSURE 最终合并后的证据列表 $R_{\text{final}}$
    \STATE $R_{\text{init}} \leftarrow \text{DenseRetrieve}(q, \mathcal{D}, k)$
    \STATE $R_{\text{scored}} \leftarrow \text{CrossEncoderRerank}(q, R_{\text{init}})$
    \STATE $s_{\text{max}} \leftarrow \max_{d \in R_{\text{scored}}} \text{score}(d)$
    \IF{$s_{\text{max}} < \tau$}
        \STATE $d_{\text{top}} \leftarrow \text{argmax}_{d \in R_{\text{scored}}} \text{score}(d)$
        \STATE $q_{\text{expand}} \leftarrow \text{ExtractKeywords}(d_{\text{top}}) \cup q$
        \STATE $R_{\text{second}} \leftarrow \text{DenseRetrieve}(q_{\text{expand}}, \mathcal{D}, k)$
        \STATE $R_{\text{final}} \leftarrow \text{ReciprocalRankFusion}(R_{\text{init}}, R_{\text{second}})$
    \ELSE
        \STATE $R_{\text{final}} \leftarrow R_{\text{scored}}$
    \ENDIF
    \RETURN $R_{\text{final}}$
    \end{algorithmic}
    \end{algorithm}

    % End of appendix
    \removeappendixsubsecmajornumbering
}